\subsection{Division Algebras}

\begin{definition} \label{def:division-algebra} \hlt{Division Algebras}\\
An associative $\mathbb{F}$-algebra $\mathcal{D}$ is a \hlt{division algebra} iff $\mathcal{D}\setminus\{0\}\subseteq\mathcal{D}^{\times}$.
\end{definition}

\begin{example} \label{ex:hamilton-quaternions} \hlt{Hamilton's Quaternions}\\
Hamilton's quaternion algebra is $\mathbb{H}\defeq\mathbb{R}1\oplus\mathbb{R}i\oplus\mathbb{R}j\oplus\mathbb{R}k$ with unit $1$,\\
relations $i^2=j^2=k^2=-1$, $ij=k=-ji$, $jk=i=-kj$, and $ki=j=-ik$, extended $\mathbb{R}$-bilinearly.\\
It is noncommutative since $ij=-ji$. For $q=r_0+r_1i+r_2j+r_3k$, define $\overline{q}\defeq r_0-r_1i-r_2j-r_3k$ and $N(q)\defeq q\overline{q}=r_0^2+r_1^2+r_2^2+r_3^2$.
If $q\neq0$, then $N(q)>0$ and $q^{-1}=N(q)^{-1}\overline{q}$, so $\mathbb{H}$ is a real division algebra.
\end{example}

\begin{figure}[H]
  \centering
  \begin{tikzpicture}[
  >=Stealth,
  every node/.style={font=\small},
  basis/.style={circle, draw=blue!60!black, fill=blue!7,
    minimum size=0.78cm, inner sep=1pt},
  fwdedge/.style={->, line width=0.9pt, draw=dartmouthgreen!70!black},
  lbl/.style={font=\scriptsize, fill=white, inner sep=1.5pt}
]
  \node[basis] (i) at (90:1.55) {$i$};
  \node[basis] (j) at (210:1.55) {$j$};
  \node[basis] (k) at (330:1.55) {$k$};

  \draw[fwdedge] (i) to[bend left=18] node[lbl, left=2pt] {$ij=k$} (j);
  \draw[fwdedge] (j) to[bend left=18] node[lbl, below=3pt] {$jk=i$} (k);
  \draw[fwdedge] (k) to[bend left=18] node[lbl, right=2pt] {$ki=j$} (i);
\end{tikzpicture}

  \caption{Quaternion basis products: the cycle $i\to j\to k\to i$ gives
  $ij=k$, $jk=i$, and $ki=j$; reversed order negates the product. Also
  $i^2=j^2=k^2=-1$, and $1$ is the unit.}
  \label{fig:quaternion-multiplication-cycle}
\end{figure}

\begin{theorem} \label{thm:division-algebra-algebraically-closed-field} \hlt{Division Algebras over Algebraically Closed Fields}\\
If $\dim_{\mathbb{F}}\mathcal{D}<\infty$, $\mathcal{D}$ is a division $\mathbb{F}$-algebra, and $\mathbb{F}$ is algebraically closed, then $\mathcal{D}=\mathbb{F}1_{\mathcal{D}}\cong\mathbb{F}$.
\begin{proof}
Let $a\in\mathcal{D}$. By Theorem \ref{thm:finite-dimensional-algebras-algebraic-elements}, $a$ is algebraic with minimal polynomial $m_a$.\\
If $m_a=pq$ with $0<\deg p,\deg q<\deg m_a$, then $0=m_a(a)=p(a)q(a)$, while minimality gives $p(a),q(a)\neq0$; this contradicts $\mathcal{D}\setminus\{0\}\subseteq\mathcal{D}^{\times}$. Thus $m_a$ is irreducible, hence linear over algebraically closed $\mathbb{F}$: $m_a=x-r$.\\
Therefore $a=r1_{\mathcal{D}}$, so $\mathcal{D}\subseteq\mathbb{F}1_{\mathcal{D}}$, and the reverse inclusion is the scalar copy.\\
\hlt{Structural Commentary}: Algebraic closure leaves only scalars.\\
\hlt{Key Insight}: Irreducible minimal polynomials collapse to linear polynomials.
\end{proof}
\end{theorem}

\begin{theorem} \label{thm:wedderburn-finite-division-algebra} \hlt{Wedderburn's Theorem}\\
Every finite division algebra $\mathcal{D}$ is commutative; equivalently $\mathcal{D}=Z(\mathcal{D})$, hence $\mathcal{D}$ is a finite field.
\begin{proof}
Let $F\defeq Z(\mathcal{D})$, $|F|=z$, $[\mathcal{D}:F]=n$, and $G\defeq\mathcal{D}^{\times}$.\\
Suppose $\mathcal{D}$ is noncommutative, and let $B$ be representatives of the noncentral conjugacy classes of $G$.\\
For $\beta\in B$, $C_{\mathcal{D}}(\beta)$ is a proper division subalgebra containing $F$, so $|C_{\mathcal{D}}(\beta)|=z^{k_\beta}$ with $1<k_\beta<n$ and $k_\beta\mid n$.\\
By the class equation, Theorem \ref{thm:finite-group-class-equation}, $z^n-1=z-1+\sum_{\beta\in B}(z^n-1)/(z^{k_\beta}-1)$.\\
Since $k_\beta\mid n$ and $k_\beta<n$, the primitive $n$th roots occur in $(x^n-1)/(x^{k_\beta}-1)$, so Propositions \ref{prop:roots-of-unity-cyclic-and-cyclotomic-factorization} and \ref{prop:cyclotomic-polynomials-integer-coefficients} give $Q_n(z)\mid(z^n-1)/(z^{k_\beta}-1)$ for each $\beta$, and also $Q_n(z)\mid z^n-1$.\\
The class equation therefore gives $Q_n(z)\mid z-1$, contradicting Lemma \ref{lem:cyclotomic-size-estimate-for-wedderburn}, which gives $|Q_n(z)|>z-1$.\\
Thus $G$ has no noncentral conjugacy classes, so $G$ is abelian and $\mathcal{D}$ is commutative.\\
\hlt{Structural Commentary}: The class equation converts noncommutativity into cyclotomic divisibility.\\
\hlt{Key Insight}: The cyclotomic lower bound is too large to divide $z-1$.
\end{proof}
\end{theorem}

\begin{theorem} \label{thm:frobenius-real-division-algebras} \hlt{Frobenius Classification over $\mathbb{R}$}\\
If $\dim_{\mathbb{R}}\mathcal{D}<\infty$ and $\mathcal{D}$ is a real division algebra, then $\mathcal{D}\cong\mathbb{R}$, $\mathcal{D}\cong\mathbb{C}$, or $\mathcal{D}\cong\mathbb{H}$.
\begin{proof}
For each $d\in\mathcal{D}$, $m_d$ is linear or irreducible quadratic over $\mathbb{R}$; completing the square gives $d=r+u$ with $r\in\mathbb{R}$ and $u^2\in\mathbb{R}_{\le0}$. Set $\mathcal{D}_0\defeq\{u\in\mathcal{D}:u^2\in\mathbb{R}_{\le0}\}$; the preceding decomposition and uniqueness give $\mathcal{D}=\mathbb{R}\oplus\mathcal{D}_0$. If $\mathcal{D}_0=0$, then $\mathcal{D}\cong\mathbb{R}$. If $\dim\mathcal{D}_0=1$, choose $i\in\mathcal{D}_0$ with $i^2=-1$, so $\mathcal{D}=\mathbb{R}\oplus\mathbb{R}i\cong\mathbb{C}$.\\
If $\dim\mathcal{D}_0>1$, use $\langle u,v\rangle\defeq-\frac12(uv+vu)$ on $\mathcal{D}_0$ to choose orthonormal $i,j\in\mathcal{D}_0$ with $i^2=j^2=-1$; $ij+ji=0$.\\
Then $k\defeq ij$ satisfies $k^2=-1$, $ij=k$, $jk=i$, $ki=j$, and reversed products change sign.\\
If $\ell\in\mathcal{D}_0$ is orthogonal to $i,j,k$, the same identities force $\ell=0$, so $\mathcal{D}=\mathbb{R}\oplus\mathbb{R}i\oplus\mathbb{R}j\oplus\mathbb{R}k\cong\mathbb{H}$.\\
\hlt{Structural Commentary}: A real division algebra splits into scalars plus a Euclidean pure part.\\
\hlt{Key Insight}: Orthogonal imaginary units generate exactly $\mathbb{C}$ or $\mathbb{H}$.
\end{proof}
\end{theorem}
